\documentclass[12pt]{article}
\usepackage[utf8]{inputenc}
\usepackage[T1]{fontenc}
\usepackage{hyperref}
\title{Gaam Language Resources for K-12 Curriculum Development}
\author{Compiled by Claude Sonnet 5 (claude-sonnet-5) for Andr\'as Kornai}
\date{}
\begin{document}
\maketitle
\section{Introduction}
Gaam (Gaahmg, Ingessana; ISO 639-3 tbi; Glottocode gaam1241; Eastern Jebel) is spoken in the Tabi Hills of Blue Nile State, Sudan. Speaker estimates vary (about 67,000 in the survey reference; Wikipedia lists about 110,000). Little educational material was found.
\section{Primary Educational Resources}
No primers, textbooks or school programmes were found in search results.
\section{Linguistic and Reference Materials}
\begin{itemize}
\item Timothy Stirtz, \emph{A grammar of Gaahmg} (Leiden doctoral dissertation, 2012). \url{https://scholarlypublications.universiteitleiden.nl/access/item%3A2963286/view}
\item Bender and Ayre, \emph{Preliminary Gaam-English-Gaam Dictionary} (1980) and a later Gaam--English dictionary (Madal, Khartoum, 2004) are named on the Wikipedia page. \url{https://en.wikipedia.org/wiki/Gaam_language}
\item Omniglot notes on Gaam writing (Latin script; tone usually unmarked). \url{https://omniglot.com/writing/gaam.htm}
\end{itemize}
\section{Religious and Community Resources}
No Bible translation, scripture portions or community organisation was found in search results.
\section{Digital and Audio Resources}
\begin{itemize}
\item Wikitongues page for tbi. \url{https://wikitongues.org/languages/tbi/}
\item Omniglot alphabet chart. \url{https://www.omniglot.com/charts/gaam.pdf}
\end{itemize}
\section{Curriculum Development Context}
A curriculum would have to start from the Stirtz grammar and the existing dictionaries, since no classroom materials were found. Orthography decisions (tone marking) will need community input.
\end{document}